\chapter{Preliminares metodologicos}

\section{Objeto del manual}
La investigacion de operaciones estudia la construccion, analisis y solucion de
modelos cuantitativos para apoyar decisiones bajo restricciones. Este manual
adopta una perspectiva doble: presenta los fundamentos clasicos del campo y
muestra como convertirlos en modelos ejecutables con \qsb.

\section{Ciclo de modelacion}
Un estudio riguroso debe registrar, como minimo:
\begin{enumerate}
  \item el sistema real y la decision que se desea apoyar;
  \item las variables de decision, parametros y unidades;
  \item los supuestos que vuelven tratable el modelo;
  \item la funcion objetivo y las restricciones;
  \item el procedimiento de solucion;
  \item la interpretacion, validacion y sensibilidad del resultado.
\end{enumerate}

\section{Trazabilidad computacional en QSB}
Los ejemplos del manual usaran archivos de entrada en
\texttt{manual/examples/} y salidas obtenidas con comandos de \qsb. Cuando un
ejemplo provenga de datos
externos, debera documentarse su fuente, licencia, fecha de extraccion,
transformaciones y unidad de analisis.

\begin{observacion}
Un resultado computacional no sustituye la interpretacion. El solucionador
produce valores optimos bajo los supuestos codificados; el analista debe decidir
si esos supuestos son defendibles para el sistema real.
\end{observacion}
